% Source-aligned English for printed p.137.
\SourcePage{137}

\LayerTarget{tohoku:thm:1.10.1:construction:transfinite:p2}
\noindent set $M_i(A)=M_1(M_j(A))$, and let the morphism
$M_j(A)\longrightarrow M_{j+1}(A)$ be $f(M_j(A))$; this simultaneously
defines the morphisms $M_k(A)\longrightarrow M_i(A)$ for $k\leq i$. If $i$ is
a limit ordinal, set

\[
 M_i(A)=\mathop{\rm lim}\limits_{\longrightarrow,\,j<i}M_j(A),
\]

\noindent and take as the morphisms $M_j(A)\longrightarrow M_i(A)$ for $j<i$
the canonical morphisms, which are indeed injective by Proposition 1.8. Let
$k$ now be the least ordinal whose cardinality is strictly greater than the
cardinality of the set of subobjects of $U$, and set $M(A)=M_k(A)$. It remains
to prove that $M_k$ is injective, that is, that it satisfies the condition of
Lemma 1. With the notation of that lemma, we prove that $v(V)$ is contained in
one of the $M_i$ with $i<k$, which will finish the proof. Indeed, from
$M_k=\mathop{\rm Sup}M_i$ one obtains

\[
 V=\mathop{\rm Sup}_{i<k}v^{-1}(M_i)
\]

\noindent by AB 5). The set of subobjects of $V$ has cardinality less than
$\mathrm{card}\,k$, while every set of ordinals below $k$ and cofinal in $k$
has cardinality $k$, since $\mathrm{card}\,k$ is not a limit cardinal.
Consequently $v^{-1}(M_i)$ is constant from some $i_0<k$ onward, so
$V=v^{-1}(M_{i_0})$. This completes the proof.

\LayerTarget{tohoku:thm:1.10.1:remark:1}
\noindent\textsc{Remarks. 1. Variant of Theorem 1.10.1:} If $C$ satisfies
axioms AB 3), AB 4), and AB 3*), and has a cogenerator $T$, then every
$A\in C$ has an injective effacement. We shall not use this result.

\LayerTarget{tohoku:thm:1.10.1:remark:2}
\noindent 2. The fact that $M(A)$ is functorial in $A$ can be useful, for
example in proving that every object $A\in C^G$ (that is, an object $A\in C$
with operator group $G$; cf. 1.7, Example f) that is injective in $C^G$ is
also injective in $C$.

\LayerTarget{tohoku:thm:1.10.1:remark:3}
\noindent 3. In many cases the existence of a monomorphism from $A$ into an
injective object can be seen directly in a simpler way. The advantage of
Theorem 1 is that it applies to very different cases. Moreover, its conditions
are stable under passage to certain categories of diagrams (cf.
Propositions 1.6.1 and 1.9.2), where the existence of enough injective objects
is not always visible to the naked eye.

\LayerTarget{tohoku:thm:1.10.1:remark:4}
\noindent 4. We leave to the reader the task of giving the dual statements for
projective objects and projective effacements.

\LayerTarget{tohoku:section:1.11}
\noindent\textbf{1.11. Quotient categories.} Although they will not be used in
the rest of this work, the considerations of this section, which systematize
and make more flexible Serre's ``language modulo $C$'' [17], are useful in
various applications.

\LayerTarget{tohoku:def:subcategory:p1}
Let $C$ be a category. A \textit{subcategory of $C$} is a category $C'$ whose
objects are objects of $C$, such that for $A,B\in C'$ the set
$\mathrm{Hom}_{C'}(A,B)$ of morphisms from $A$ to $B$ in the sense of $C'$ is
a subset of the set $\mathrm{Hom}_C(A,B)$ of morphisms from $A$ to $B$ in the
sense of $C$, composition in $C'$ is induced by composition in $C$, and the
identity morphisms in $C'$ are identity morphisms in $C$. The last two
conditions say that the ``function'' assigning to an object or morphism of
$C'$ the same object or morphism of $C$ is a covariant functor from $C'$ to
$C$, called the \textit{canonical injection of $C'$ into $C$}. If $C$ and
$C'$ are additive categories, then $C'$ is called an \textit{additive
subcategory} if

\TohokuContinuation{138}
